\documentclass[12pt]{article}
\usepackage[utf8]{inputenc}
\usepackage[T1]{fontenc}
\usepackage[margin=1in]{geometry}
\usepackage{setspace}
\usepackage{xcolor}
\usepackage[colorlinks=true, urlcolor=blue!60!black, linkcolor=blue!60!black]{hyperref}
\setlength{\parindent}{0em}
\setlength{\parskip}{0.6em}
\renewcommand{\baselinestretch}{1.0}
\begin{document}
\begin{center}
{\LARGE \textit{Sherbert v.\ Verner}, 374 U.S.\ 398 (1963)\par}
\vspace{0.5em}
{\large Case Brief --- Two-Page Edition\par}
\end{center}

\section{Memory Jogger}
Denying unemployment benefits to a Sabbatarian who would not work Saturdays violated free exercise; burdens on religion need a ``compelling state interest'' (\href{https://supreme.justia.com/cases/federal/us/374/398/}{Justia}).

\section{Detailed Case Facts}
Adell Sherbert, a mill worker for about 35 years, became a Seventh-day Adventist. When her mill moved to a mandatory six-day week she refused Saturday work and was fired, and she could find no five-day job. South Carolina denied benefits for failing ``without good cause'' to accept suitable work.

\section{Procedural History}
The Employment Security Commission denied benefits; the county court and the South Carolina Supreme Court affirmed. The U.S.\ Supreme Court reversed.

\section{Judicial Votes}
Reversed 7--2. Brennan wrote for Warren, Black, Douglas, Clark, and Goldberg; Douglas concurred; Stewart concurred in the result; Harlan (with White) dissented.

\section{Holding}
A State may not force a worker to choose between her faith and public benefits unless a compelling interest justifies the burden and no less restrictive alternative exists. Extending benefits does not establish religion.

\section{Analysis of Opinions}
\textbf{Majority:} the choice works like a fine on Saturday worship; benefits need not be a ``right''; the State's fear of fraud was speculative; and \textit{Braunfeld} was distinguishable. \textbf{Stewart:} the decision collides with ``positively wooden'' Establishment Clause cases and \textit{Braunfeld}. \textbf{Harlan:} the neutral rule disqualified all who were unavailable for personal reasons; accommodation is permitted, not compelled.

\section{Examples: Future Applications}
\textbf{Same-Side:} a pacifist who quits rather than build tanks keeps benefits (\textit{Thomas v.\ Review Board}). \textbf{Opposite-Side:} a neutral peyote ban applied to religious use (\textit{Employment Division v.\ Smith}). \textbf{Fence-Sitter:} a worker refusing Saturday shifts for a secular moral duty.

\section{Critique}
Originalists debate whether the Founders contemplated judicial exemptions; critics say \textit{Sherbert} invites judges to weigh religious claims. Equality scholars warn that privileging religion over secular conscience undermines equal citizenship.

\section{Key Quotations}
\begin{quote}
``The ruling forces her to choose between following the precepts of her religion and forfeiting benefits, on the one hand, and abandoning one of the precepts of her religion in order to accept work, on the other hand. Governmental imposition of such a choice puts the same kind of burden upon the free exercise of religion as would a fine imposed against appellant for her Saturday worship.''
\end{quote}
\noindent --- Justice Brennan, for the Court, 374 U.S.\ at 404.

\end{document}
